% This LaTeX file uses the class impletter.cls to format
% the new Imperial College letter to be printed on blank paper.
% LdC Foulkes Jan. 2003.

\documentclass[blank]{impletter}
\usepackage{graphicx}

%%%%%%%%%%%%%%%% Adjustment Parameters %%%%%%%%%%%%%%%%%%%%%%%%%%%
% The letter has to have have the following (measured from the top):
%	The bottom of the word 'Imperial' in the logo at 15mm
%	The bottom of the date at 51mm
%	The bottom of the 1st address line at 63mm
%	The bottom of the salutation at 105mm
%	The bottom of the department line (top right) at 12mm
%	The bottom of the statutory footer 5.5mm from the bottom
% 	The bottom of the first line of text on the second page at 51mm
%	The margin at the left of the text should be 25mm
%	The sender's detail (top right block) should be 90mm from the
%	right edge.

% The following are adjustable parameters to account for printer
% differences. The values can be given in mm or pt
% (a point is about .35 of a mm).

% To move the printed page up (minus value) or down (plus value):
\addtolength{\topmargin}{0mm}

% To move the printed text to the right (plus) or left (minus).
\addtolength{\evensidemargin}{0pt}
\addtolength{\oddsidemargin}{0pt}

% To move the statutory footer up (minus) or down (plus) - useful
% if the printer will not print so close to the bottom of the page
\addtolength{\footskip}{0mm}
%%%%%%%%%%%%%%%%%%%%%%%%%%%%%%%%%%%%%%%%%%%%%%%%%%%%%%%%%%%%%%%%%%

\begin{document}

\headers{
% Replace with the addressee details but DO NOT remove the \\ even if empty
NSF GRFP Review Committee
}
% Replace with the salutation for this letter
{
Dear Committee,
}
% Replace with your department, address, telephone, fax and e-mail.
% Do not remove the blank lines.
{Department of Life Sciences\\
Imperial College London, Silwood Park Campus\\
Buckhurst Road, Ascot\\
Berkshire SL5 7PY, UK\\
Telephone +44 (0)20 7594 2213\\
s.pawar@imperial.ac.uk\\
www.imperial.ac.uk/people/s.pawar\\
}
% Replace with your name
{
Dr Samraat Pawar
}
%% Replace with your qualifications
% {
% BA, MS, PhD, CPhys, MinstP\\
% }
% Replace with your job title. A second line can be added if necessary.
{
Lecturer\\
Grand Challenges in Ecosystems and the Environment Initiative\\
% Director, MSc/MRes in Computational Methods in Ecology\\
% \& Evolution 
}
\informal

% The text of the letter starts here

I am writing this letter in support of Matthew Smith's (Matt's) application to 
the NSF GRFP. 

I have known Matt since July 2012. I was a postdoctoral scholar in Dr. Stefano 
Allesina's lab then, and Matt had just joined as a post-baccalaureate 
researcher. At the very outset, I was impressed by Matt's work ethic his ability 
to efficiently and effectively learned new and difficult techniques, both at the 
computational and mathematical end of things. For example, within a few months, 
he went from somebody who was beginning to learn the R programming language to a 
veritable reference on its syntax and usage! Over the year that we shared lab 
space, I had many discussions with Matt about various conceptual and theoretical 
issues in ecology and evolutionary biology, and I was always impressed by his
clarity of 
thinking and knack of approaching problems in unconventional ways.

When I left the University of Chicago in July 2013 to take up my current 
position, Matt has just joined the PhD program in Ecology and Evolution at the 
University of Chicago. Based upon my interactions with previous years'
graduate 
student cohorts, I would easily place Matt among the top 5\% of the 
Department's Graduate students. Clearly, my assessment is well backed by Matt's 
outstanding academic and extracurricular record in his High School and BS 
years. 

I will now make some comments about the merit of Matt's proposed research 
project, which is strongly related the research focus of my lab group at 
Imperial College. But first, I would like to emphasize that both the areas that 
Matt worked on during his pre-PhD work at U Chicago --- network-topographical 
model selection and localization of eigenvalues of linear operators (e.g., 
the Jacobian matrix) are not only important, but also technically challenging 
areas of research. I literally shared desk space with Matt while he was working 
on these projects, and had many discussions with him about them. I was 
consistently impressed by his ability to explain the technical details
of these 
projects, and by his persistence and ability in mastering the necessary 
techniques. Just as importantly, I was impressed by his ability and motivation 
to seek where and how both these topics of research could be applied to real 
world problems.

The work that Matt is now proposing, which is a continuation of one of his two 
Undergraduate research projects in the Allesina lab, is an exciting and 
important avenue of research. The study of dynamics and stability of biological 
systems is currently a rapidly growing area of research. This is not surprising, 
given that we live in an environment that is increasingly becoming uncertain, 
and biological systems, not the least ecological networks such as food webs, are 
under considerable threat. Although ecological networks are complex, nonlinear 
systems, the study of their behavior in the neighborhood of local, stable points 
is extremely important --- any point that is locally stable also has
the 
potential to be non-locally (or even globally) stable. The main focus of Matt's 
project --- the localization of eigenvalues of the Jacobian Matrix of 
finitely-sized and incompletely connected matrices --- is extremely important in 
real world applications, because, well, real world systems are finitely-sized 
and incompletely connected! His proposed research is all the more important 
because it will directly address the fact that to  great extent, this 
incompleteness is a result of inadequate sampling, an unavoidable problem in the 
real world. Moreover, I think the approach of combining moment-closure and 
Maximum entropy that Matt has taken is imaginative and promising, and is highly 
likely to yield new and general results. 

Finally, I would like to emphasize that Matt has incredible awareness of 
and sensitivity to current social and ecological problems. In fact, in my/ 
opinion, Matt's path from field-oriented to theoretical research has placed him 
in a rare position where he can always keep real-world applications in mind 
while conceiving and carrying out his theoretical research. I am confident 
that he will take exceptional effort and will succeed at extending his work to 
real world applications, and disseminating his research on public as well as 
academic platforms. His current involvement in outreach and education in the 
Hyde Park community in Chicago speaks volumes of his motivation for
outreach 
and commitment to society.

So in summary, I enthusiastically and wholeheartedly support Matthew Smith's 
application to the NSF GRFP. This fellowship award will aid in the making of an 
exceptional research and academic career.

% Leave the name field blank if you don't want your name to be
% printed at the bottom, but do not remove the {}.
\close{
Yours sincerely,
\vspace{-40pt}
}
{
\includegraphics[width=0.3\linewidth]{graphics/samsig.pdf}\\
Samraat Pawar
}
\end{document}
